%****************************************************
%Choose your desired class
%
%\documentclass[10pt, a4paper, two-sided,]{scrreport}
\documentclass[10pt, a4paper, two-sided,]{scrbook}
%\documentclass[a4paper]{letter}
%\documentclass[a4paper]{article}
%
%****************************************************

\usepackage{lipsum} 

\usepackage[swedish,english]{babel}
\usepackage[T1]{fontenc}
\usepackage[utf8]{inputenc}
\usepackage{eso-pic}
\usepackage{graphicx}           % Prepare for adding images
\usepackage{hyperref}           % Get refs, ToC, etc hyperlinked

\setlength{\parindent}{0pt}
\pagestyle{empty}               % No page number on title page and abstract page

%\setlength{\headsep}{5cm}
%\setlength{\oddsidemargin}{5mm}

\newcommand{\thesistitle}[2]{\textsf{\Huge #1}\par\vspace{5mm}\textsf{#2}\par\vspace{10mm}}
\newcommand{\thesisauthor}[1]{\textsf{\Large #1}}
\newcommand{\supervisor}[1]{\textsf{\Large Handledare: #1}}
\newcommand{\examiner}[1]{\textsf{\Large Examinator: #1}}
\newcommand{\submissiondate}[1]{\textsf{\Large Inlämningsdatum: #1}}
\newcommand{\Framsida}{\AddToShipoutPicture*{\put(0,0){\includegraphics{kandidatfram.pdf}}}}
\newcommand{\Baksida}{\backmatter\pagestyle{empty}\AddToShipoutPicture*{\put(0,0){\includegraphics{kandidatbak.pdf}}}}



\begin{document}

\Framsida %Kommentera för att bli av med loggan
\vspace*{4cm}
\thesistitle{Fin titel}{Nice title} %Fyll i titel. Använd först uppsatsens språk och sen det "andra" språket (oftast svenska)
\thesisauthor{Augusta Ada King-Noel, Countess of Lovelace} %Fyll i författare

% \vspace*{11cm}
\vfill{}
\supervisor{Marty McFly} \medskip
 
\examiner{Emmett Brown} \medskip

\submissiondate{2015-10-21} \medskip

\pagebreak{}
\begin{center}
	\textbf{Abstract}
\end{center}
\textit{[Place your English abstract here]}
\lipsum[1]

\bigskip
\begin{center}
	\textbf{Sammanfattning}
\end{center}
\textit{[Placera ditt svenska abstrakt här ]}
\lipsum[1]

\newpage

% HERE the content of your thesis incl.\ table of contents, main body, bibliography etc
\tableofcontents{}              % Remove or comment out if you don't want a table of contents


\pagestyle{plain}               % We want pagenumbering from here

\section*{Some LaTeX advice}
\label{sec:some-latex-advice}

\begin{itemize}
\item Use the \emph{enumerate} or \emph{itemize} environments when you want lists of some kind. The
  package \emph{enumitem} can be used to give you more control over spacing and details of these
  lists.
\item It is considered better style to write ``like this'' than "like that" in LaTeX. This does not
  seem to be well-supported in Overleaf, unfortunately.
\end{itemize}

\section*{Advice on figures}
\label{sec:advice-figures}

\begin{itemize}
\item Put figures in a \emph{figure} environment and make sure to have a label (like ``fig:accuracy'') to refer to. 
\item Avoid referring to figures with ``see below'' or similar. Use \verb+see Figure~\ref{fig:accuracy}+ and similar instead. The tilde is a hard space that avoids line breaks.
\item Do not ``hard code'' numbers for figures, tables, et.c. At one point or another you will forget to update the numbers. LaTeX gives you great support for automatic numbering with \verb|\label{somelabel}| and \verb|\ref{somelabel}|.
\item Make sure you have the rights to add figures that you do not make yourself. You thesis is
  public, so published, and falls under copyright laws. It is easy to find figures/pictures that you
  may use, but you need to mention the author/source.
\end{itemize}



\chapter{Introduction}
\label{cha:introduction}

Give a general background to the project.


\section{Some section}
\label{sec:some-section}

You might want to structure the intro somewhat.




\lipsum[1-3]

\chapter{Methods}
\label{cha:methods}

\lipsum[4-5]


%
% Possibly an appendix? Often not needed, but sometimes useful. This is marked with a special section.
% Simply remove if you don't see a use for this.
% 
\appendix

\chapter{Extra stuff}
\label{cha:extra-stuff}

\section{Perhaps extra figures}
\label{sec:perh-extra-figur}


\lipsum[6]

\section{Perhaps extra tables}
\label{sec:perhaps-extra-tables}



\lipsum[7]

\Baksida %Kommentera för att bli av med loggan
\vspace*{13cm}
%Kandidatuppsats 2010:11\\ %Fyll i löpnummer
Datalogi\\
%September 2010\\\\ % Fyll i månad och år
www.math.su.se\\\\
Beräkningsmatematik\\
Matematiska institutionen\\
Stockholms universitet\\
106 91 Stockholm\\





\end{document}

%%% Local Variables:
%%% mode: latex
%%% TeX-master: t
%%% End:
